\section*{第 \ref{ch:b3:sensory-systems} 章——感觉系统}

\begin{solution}{exo:b3:sensory-systems:1}
由线路而非信息决定：每根感觉神经都是一条\omterm{def:b3:sensory-systems:principles}{标记线路}，
终止于大脑中各自专属的区域；凡沿视神经传来的，都被读作光，
所以压迫眼球会产生闪光，击打耳部则
产生耳鸣。
\end{solution}

\begin{solution}{exo:b3:sensory-systems:2}
暗处时，环 GMP 维持阳离子通道开放，
内向的稳定“暗电流”使细胞保持去极化并释放谷氨酸。
光激活\omterm{def:b3:sensory-systems:retina}{视紫红质}、\omterm{def:b3:bacteriology:hgt}{转导}蛋白和磷酸二酯酶，后者分解环 GMP；
通道关闭，内向电流停止，
细胞超极化，递质释放减少——
光以信号减弱的方式被传递。
\end{solution}

\begin{solution}{exo:b3:sensory-systems:3}
$I = I_{0}10^{L/10}$：\qty{20}{dB}，\qty{1e-10}{W/m^2}；\qty{60}{dB}，
\qty{1e-6}{W/m^2}；\qty{110}{dB}，\qty{0.1}{W/m^2}。由响到轻：
$10^{9}$。
\end{solution}

\begin{solution}{exo:b3:sensory-systems:4}
甜、鲜和苦味经由 G 蛋白偶联受体感知（甜味、鲜味受体各一种，
苦味受体约二十五种）；咸味经钠通道 ENaC 感知，
酸味经质子通道感知。“风味”的大部分是挥发性分子
由口腔经鼻后通道进入鼻腔、被嗅到的气味；鼻子被堵住时，
只剩五种\omterm{def:b3:sensory-systems:chemical}{味觉}，食物便寡淡无味。
\end{solution}

\begin{solution}{exo:b3:sensory-systems:5}
费希纳：$S = (1/0.02)\ln(10\,000/100) = 50\ln 100 \approx 230$ 级。
计数法：$1.02^{n} = 100$ 得 $n = \ln 100/\ln 1.02 = 233$。
\end{solution}

\begin{solution}{exo:b3:sensory-systems:6}
$P = 1 - e^{-a}\sum_{k=0}^{5}a^{k}/k!$：$a = 2$：$0.017$；$a = 6$：
$0.55$；$a = 12$：$0.98$。当 $a \approx 5.7$ 时，半数闪光可被看到；
若吸收率为 \qty{10}{\%}，则角膜处须约有 $57$
个光子。
\end{solution}

\begin{solution}{exo:b3:sensory-systems:7}
$55/3\times 1.3 \approx 24$；以声压\omterm{def:b3:sensory-systems:cochlea}{分贝}计，$20\log_{10}24
\approx \qty{28}{dB}$。听小骨融合，无法传导振动，因此
声音只能经颅骨传导抵达\omterm{def:b3:sensory-systems:cochlea}{耳蜗}，
无中耳增益，即产生
约 \qtyrange{30}{50}{dB} 的传导性听力损失。
\end{solution}

\begin{solution}{exo:b3:sensory-systems:8}
\omterm{def:b3:sensory-systems:cochlea}{毛细胞}的通道由\omterm{def:b3:sensory-systems:cochlea}{顶连接}直接牵拉打开，
无需信使，亦无需酶，故响应只需微秒量级；光感受器则
借两个酶促阶段放大单个光子，这两级各需数十毫秒。听觉
获得了借双耳延迟定位声源、追踪相位所需的时间精度；
视觉则获得了计数单个光子的灵敏度，
代价是速度变慢。
\end{solution}

\begin{solution}{exo:b3:sensory-systems:9}
$H(0.05) = -0.05\ln 0.05 - 0.95\ln 0.95 = 0.150 + 0.049 = 0.199$。
人类：$\ln\binom{400}{20} \approx 400\times 0.199 = 80$，约
$10^{34}$ 种模式。小鼠：$\ln\binom{1100}{55} \approx 1100\times 0.199
= 219$，约 $10^{95}$。
\end{solution}

\begin{solution}{exo:b3:sensory-systems:10}
每时间窗的平均自发事件数为 $500\times 0.1/40 = 1.25$。$P(k
\le 5) = e^{-1.25}(1 + 1.25 + 0.78 + 0.33 + 0.10 + 0.03) = 0.998$，故
$P(k \ge 6) \approx 2\times 10^{-3}$。若阈值取一，则
一次自发事件将出现在多数时间窗中（$1 - e^{-1.25} = 0.71$），
而黑暗中将满是闪光；取六时，误报每
五百个时间窗才出现一次。阈值应由
\omterm{def:b3:sensory-systems:retina}{视杆细胞}的噪声而非其敏感度决定。
\end{solution}

\begin{solution}{exo:b3:sensory-systems:11}
最大延迟 $0.2/340 = \qty{590}{\micro s}$（来自一侧的声音）。
在中线附近 $\Delta t \approx (d/c)\sin\theta$，故 $\qty{10}{\micro
s}$ 对应 $\sin\theta = 0.017$，约 \qty{1}{\degree}。高于
\qty{4}{kHz} 时，周期（\qty{250}{\micro s}）短于
神经元能够锁定的时间尺度，而数百微秒的延迟
跨越了一个以上的周期，使相位变得含糊；此时大脑转而利用
两耳间的强度差与位置编码。
\end{solution}

\begin{solution}{exo:b3:sensory-systems:12}
远离边缘时，$r = I(1 - 2w)$：暗侧为 $0.6$，亮侧为
$1.2$。在最后一个暗侧感受器处（$I = 1$，相邻为 $1$ 和 $2$）：$r
= 1 - 0.2\times 3 = 0.4$；在第一个亮侧感受器处：$r = 2 - 0.6 = 1.4$。
阶跃被放大为一个负向过冲和一个正向过冲——即马赫
带。\omterm{def:b3:sensory-systems:retina}{视网膜}由此获得对比度与边缘，并以较小的动态范围
进行传递；但丧失了对绝对亮度的\omterm{prop:b3:dna-repair:fidelity}{保真度}，并产生
亮度错觉。
\end{solution}

\begin{solution}{pb:b3:sensory-systems:1}
\textbf{1.} $0.01/3.6\times 10^{-19} = 2.8\times 10^{16}$ 个光子每
秒。
\textbf{2.} 瞳孔面积 $\pi(3.5\times 10^{-3})^{2} = 3.8\times 10^{-5}$
m$^{2}$。在 \qty{1}{km} 处，比例为 $3.8\times 10^{-5}/1.26\times 10^{7} =
3\times 10^{-12}$，$8.5\times 10^{4}$ 个光子每秒。在 \qty{10}{km} 处：
$850$ 个光子每秒。
\textbf{3.} 在 \qty{0.1}{s} 内，吸收率为 \qty{10}{\%}：$850$ 于
\qty{1}{km} 处，$8.5$ 于 \qty{10}{km} 处——两者均高于六；烛光
在十千米处仍可被看到（大部分时间如此）。
\textbf{4.} 每时间窗吸收六个光子，需每秒 $600$ 个光子进入
瞳孔：$4\pi d^{2} = 3.8\times 10^{-5}\times 2.8\times 10^{16}/600
= 1.8\times 10^{9}$ m$^{2}$，$d \approx \qty{12}{km}$。此时，$P(k \ge
6 \mid a = 6) = 0.55$：烛光约在半数时间窗中可被看到。
\textbf{5.} $e^{-6} = 0.0025$。阈值处的计数在各时间窗之间
波动——$6 \pm 2.4$——故光源似乎时隐时现：
即暗星的闪烁现象（大气还会附加其自身的
闪烁）。
\textbf{6.} 背景的涨落，即 $\sqrt{10^{4}} = 100$ 个光子
每时间窗，淹没了烛光的 $8.5$：探测要求信号
超过背景噪声，而非暗适应眼的阈值。
\textbf{7.} $\qty{1}{pA}/200 = \qty{5}{fA}$，每个通道；$5\times
10^{-15}/1.6\times 10^{-19} = 3\times 10^{4}$ 个离子每秒。
\textbf{8.} 三十分之一；约三十个同时到达的光子可关闭每
个通道，使\omterm{def:b3:sensory-systems:retina}{视杆细胞}饱和。
\textbf{9.} $10^{-12}\times 0.2 = 2\times 10^{-13}$ C，$1.3\times 10^{6}$ 个
阳离子：一个光子可控制超过一百万个电荷。
\textbf{10.} \omterm{def:b3:sensory-systems:retina}{视杆细胞}在数毫秒内即饱和，所有通道均已关闭，
其\omterm{def:b3:sensory-systems:retina}{视紫红质}漂白速度快于再生；\omterm{def:b3:sensory-systems:retina}{视锥细胞}，
增益较低、关闭较快，仍可持续响应，并使其
工作范围跨越六个数量级的光强。
\textbf{11.} 高增益需要长时间积分（每一级都需耗时，且
响应需延长以便累积），故\omterm{def:b3:sensory-systems:retina}{视杆细胞}慢而敏感；
\omterm{def:b3:sensory-systems:retina}{视锥细胞}放大倍数较小、关闭较快，
\omterm{thm:b3:rna-regulation:mirna-kinetics}{响应时间}为 \qty{20}{ms}，足以跟上运动，代价是
需要一百个光子。
\textbf{12.} $10^{8}/40 = 2.5\times 10^{6}$ 次自发异构化每
秒，遍及\omterm{def:b3:sensory-systems:retina}{视网膜}——但每时间窗内，任一区域仅 $1.25$ 次，
而该区域由 $500$ 个\omterm{def:b3:sensory-systems:retina}{视杆细胞}组成，远低于六的阈值，故噪声可逐区剔除；
阈值与\omterm{def:b3:nervous-systems:circuits}{汇聚}正是为此而设。
\textbf{13.} \qty{0.1}{W/m^2}；在 \qty{5.5e-5}{m^2} 面积上，为 \qty{5.5e-6}{W}。
\textbf{14.} $5.5\times 10^{-6}\times 7200 = \qty{0.04}{J}$——约
相当于两小时内四百滴雨滴的能量。
\textbf{15.} $10^{-12}\times 5.5\times 10^{-5} = 5.5\times 10^{-17}$ W；
在 \qty{10}{ms} 内，$5.5\times 10^{-19}$ J——约相当于一个或
两个可见光光子的能量。
\textbf{16.} $0.3/5000 = 6\times 10^{-5}$ 弧度，$0.003^{\circ}$；其
偏转量约相当于三个氢原子。
\textbf{17.} $10^{(110 - 85)/10} = 316$：在 \qty{110}{dB} 下两小时
所提供的能量相当于在 \qty{85}{dB} 下 $630$ 小时，几乎
相当于一个月的工作日。
\textbf{18.} 约 $120$ 个恰可察觉的级，每\omterm{def:b3:sensory-systems:cochlea}{分贝}一级。
\textbf{19.} $2^{400} = 10^{400\log_{10}2} = 10^{120}$。
\textbf{20.} $\ln\binom{400}{20} \approx 400\,H(0.05) \approx 80$，减去
约为 $2$ 的斯特林修正：约 $10^{34}$ 种模式。
\textbf{21.} 计数在一种模式中有而另一种中无的受体：$5 +
5 = 10$，占 $40$（汉明距离）；重叠四分之三的
模式几乎驱动相同的\omterm{def:b3:sensory-systems:chemical}{小球}，
故被读作几乎相同的气味。
\textbf{22.} 仅占 $10^{12}/10^{34} = 10^{-22}$ 的模式。编码本身并非
限制：受体有噪声，许多受体调谐相似，因而它们的
活性相关，真实分子只能产生其中极小的一部分
模式，且大脑辨别和记忆模式的能力
亦是有限的。
\textbf{23.} $\ln\binom{1100}{20} \approx 1100\,H(0.018) = 1100\times
0.091 \approx 100$，约 $10^{43}$ 种——约为人类
二十受体模式数的 $10^{9}$ 倍（若小鼠每种
气味使用的受体更多，则远不止于此）。
\textbf{24.} 总体嗅觉几乎不变——编码具有冗余，
其他受体仍可承载每一种气味分子；但某一气味分子若依赖
该受体且浓度较低，便会无法察觉，或改变其
性质：即特异性嗅觉缺失，如许多人无法
嗅到雄烯酮。
\textbf{25.} 烛光可见距离约 \qty{12}{km}；阈值声音
在 \qty{10}{ms} 内向鼓膜传递 $5.5\times 10^{-19}$ J，约
相当于一两个光子；人类鼻子可形成约 $10^{34}$ 种二十受体模式。
\end{solution}
